\DocumentMetadata{
  pdfversion=2.0,
  pdfstandard=ua-2,
  lang=en-US,
  tagging=on,
  tagging-setup={math/setup=mathml-SE,tabsorder=structure}
}
\documentclass[11pt,letterpaper]{article}
\usepackage[margin=0.68in,includefoot]{geometry}
\usepackage{newcomputermodern}
\usepackage{amsmath,amscd,mathtools}
\usepackage{tikz,adjustbox}
\providecommand{\square}{\mdlgwhtsquare}
\usepackage[hidelinks]{hyperref}
\hypersetup{
  pdftitle={Algebra First-Year Exam, January 2015, Part 2},
  pdfauthor={University of Florida Department of Mathematics},
  pdfsubject={Graduate mathematics examination},
  pdfdisplaydoctitle=true,
  bookmarksopen=true
}
\setcounter{secnumdepth}{0}
\setlength{\parindent}{0pt}
\setlength{\parskip}{0.28em}
\setlength{\emergencystretch}{3em}
\setlength{\leftmargini}{2.15em}
\setlength{\leftmarginii}{2.65em}
\setlength{\leftmarginiii}{3em}
\pagestyle{empty}
\raggedbottom
\allowdisplaybreaks
\NewTaggingSocketPlug{block/list/label}{examproblem}{%
  \tagstructbegin{tag=Lbl}\tagstructend%
  \tagstructbegin{tag=\LBody}%
  \tagstructbegin{tag=\ProblemHeadingTag,title-o={\ProblemHeadingText}}%
  \tagmcbegin{tag=\ProblemHeadingTag,actualtext={\ProblemHeadingText}}%
  #2%
  \tagmcend%
  \tagstructend%
}
\newenvironment{examproblems}[2]{%
  \def\ProblemHeadingTag{#2}%
  \AssignTaggingSocketPlug{block/list/label}{examproblem}%
  \begin{enumerate}%
  \setcounter{enumi}{#1}%
  \renewcommand{\labelenumi}{\textbf{\theenumi.}}%
  \setlength{\topsep}{0.35em}%
  \setlength{\partopsep}{0pt}%
  \setlength{\itemsep}{0.48em}%
  \setlength{\parsep}{0.18em}%
}{\end{enumerate}}
\newenvironment{examsubparts}{%
  \AssignTaggingSocketPlug{block/list/label}{default}%
  \begin{enumerate}%
  \setlength{\topsep}{0.18em}%
  \setlength{\partopsep}{0pt}%
  \setlength{\itemsep}{0.16em}%
  \setlength{\parsep}{0.1em}%
}{\end{enumerate}}
\newenvironment{examinstructions}{%
  \par\begingroup\itshape
}{%
  \par\endgroup\vspace{0.18em}
}
\makeatletter
\renewcommand\subsection{\@startsection{subsection}{2}{0pt}%
  {-1.25ex plus -.25ex minus -.1ex}%
  {0.55ex plus .1ex}%
  {\normalfont\large\bfseries}}
\renewcommand{\@maketitle}{%
  \newpage\null\vskip -1em%
  {\footnotesize\bfseries
    \href{https://www.ufl.edu/}{University of Florida}, \href{https://math.ufl.edu/}{Department of Mathematics}\par}%
  \vskip -0.5em%
  \tagstructbegin{tag=H1,title={Algebra First-Year Exam, January 2015, Part 2}}%
  \tagpdfparaOff%
  \tagmcbegin{tag=H1}%
    {\LARGE\bfseries\href{https://gma.math.ufl.edu/exams/algebra-first-year/}{Algebra First-Year Exam}, January 2015, Part 2\par}%
  \tagmcend%
  \tagpdfparaOn%
  \tagstructend%
  \vskip 0.25em%
  \tagmcbegin{artifact}%
    \hrule height 1pt%
  \tagmcend%
  \vskip 1em%
}
\makeatother
\title{Algebra First-Year Exam, January 2015, Part 2}
\author{}
\date{}
\begin{document}
\maketitle
\thispagestyle{empty}
\begin{examinstructions}
Answer four problems. If you turn in more than four, only the first four will be graded. Within reason, you may use theorems as long as you state them clearly.
\end{examinstructions}
\begin{examproblems}{0}{H2}
\def\ProblemHeadingText{Problem 1.}
\item
Suppose~\(R\) is an integral domain.
\begin{examsubparts}
\item[\textnormal{(i)}]
Show that the annihilator of any finitely generated torsion~\(R\)-module is nonzero.
\item[\textnormal{(ii)}]
Give an example of a torsion~\(\mathbb{Z}\)-module whose annihilator is zero.
\end{examsubparts}
\def\ProblemHeadingText{Problem 2.}
\item
Let~\(R\) be the ring
\[
R=\{a+b\sqrt{-2}\mid a,b\in\mathbb{Z}\}\subseteq\mathbb{C}
\]
 (you may assume it is a ring).
\begin{examsubparts}
\item[\textnormal{(i)}]
Show that~\(R\) is a unique factorization domain.
\item[\textnormal{(ii)}]
Find the group of units of~\(R\).
\end{examsubparts}
\def\ProblemHeadingText{Problem 3.}
\item
Let~\(k\) be a field, \(V\) a~\(k\)-vector space of finite dimension, and \(V^*=\operatorname{Hom}(V,k)\) be the dual vector space. Show that \(\dim V=\dim V^*\).
\def\ProblemHeadingText{Problem 4.}
\item
Suppose~\(R\) is a commutative ring with identity, and \(S\subseteq R\) is a multiplicative subset, i.e. \(0\notin S\), \(1\in S\), and \(a,b\in S\) implies \(ab\in S\). Show that there is a prime ideal \(P\subset R\) such that \(P\cap S=\varnothing\).
\def\ProblemHeadingText{Problem 5.}
\item
Which of the following polynomials are irreducible in their respective rings? Explain.
\begin{examsubparts}
\item[\textnormal{(i)}]
\(X^3-X^2-X-2\in\mathbb{Q}[X]\)
\item[\textnormal{(ii)}]
\(Y^3+X^2Y+X^3Y+X\in k[X,Y]\), \(k\) a field
\end{examsubparts}
\def\ProblemHeadingText{Problem 6.}
\item
Let \(G=\mathrm{GL}_3(\mathbb{F}_2)\) be the group of \(3\times3\) invertible matrices with entries in the field with two elements.
\begin{examsubparts}
\item[\textnormal{(i)}]
Find representatives for all the conjugacy classes of elements of order~\(2\).
\item[\textnormal{(ii)}]
Do the same for elements of order~\(7\).
\end{examsubparts}
\end{examproblems}
\end{document}
